

\section{GUI} 

\subsection{Command Handling}
\begin{enumerate}
    \item \textbf{CONNECT Command}

    Syntax: \texttt{CONNECT <IP\_ADDRESS>}

    \begin{enumerate}
        \item Establishes TCP/IP connection to the remote measurement server on port 50012 using the specified IPv4 address
        \item Implements non-blocking socket creation with 3-second timeout to prevent GUI freezing during connection attempts
        \item Upon successful connection, automatically spawns a dedicated pthread for asynchronous network data reception
        \item Transitions application state from DISCONNECTED to CONNECTED and pre-selects Temperature and ADC1 sensors for immediate visualization
        \item Example: \texttt{CONNECT 192.168.1.10}
    \end{enumerate}

    \item \textbf{DISCONNECT Command}

    Syntax: \texttt{DISCONNECT}

    \begin{itemize}
        \item Gracefully terminates the active TCP connection and stops the network receive thread using \texttt{pthread\_join()} for safe cleanup
        \item Can only be executed when in CONNECTED state; returns error if attempted while streaming is active (must STOP first)
        \item Closes socket file descriptor, clears all sensor circular buffers, and resets timestamp synchronization reference (\texttt{server\_t0})
        \item Transitions application state back to DISCONNECTED, re-enabling the IP entry field and Connect button
        \item Example: \texttt{DISCONNECT}
        \item Error Example: \texttt{DISCONNECT} \\
        \(\rightarrow\) Error: ``Cannot disconnect: GUI is running, stop and disconnect.'' \\
        (when attempted during RUNNING state)
    \end{itemize}

    \item \textbf{START Command}

    Syntax: \texttt{START}

    \begin{itemize}
        \item Sends ASCII \texttt{"START\textbackslash n"} command to the server, initiating real-time binary sensor data streaming over the established TCP connection
        \item Enables sensor selection checkboxes with 2-sensor maximum constraint and activates frequency configuration controls
        \item Transitions state from CONNECTED to RUNNING; disables Disconnect and Shutdown buttons to prevent unsafe operations during streaming
        \item Server responds with initial RATES configuration message followed by continuous batches of \texttt{sensor\_data\_t} structures
        \item Example: \texttt{START}
        \item Typical Workflow Example:
        \begin{enumerate}
            \item \texttt{CONNECT 192.168.1.10}
            \item \texttt{START}
            \item (Data streaming begins, graph updates in real-time)
        \end{enumerate}
    \end{itemize}

    \item \textbf{STOP Command}

    Syntax: \texttt{STOP}

    \begin{itemize}
        \item Sends ASCII \texttt{"STOP\textbackslash n"} command to halt sensor data streaming while maintaining the TCP connection for future operations
        \item Preserves all captured data in circular buffers allowing graph inspection after streaming stops; data persists until disconnect or reconnect
        \item Transitions state from RUNNING to CONNECTED, re-enabling Disconnect and Shutdown buttons while disabling sensor configuration controls
        \item Does not terminate the network receive thread; connection remains active and ready to resume streaming with subsequent START command
        \item Example: \texttt{STOP}
        \item Typical Workflow Example:
        \begin{enumerate}
            \item \texttt{START}
            \item (Observe data for analysis)
            \item \texttt{STOP}
            \item (Graph frozen, data preserved for inspection)
            \item \texttt{START} (can restart streaming)
        \end{enumerate}
    \end{itemize}

    \item \textbf{SHUTDOWN Command}

    Syntax: \texttt{SHUTDOWN}

    \begin{itemize}
        \item Displays modal confirmation dialog before sending \texttt{"SHUTDOWN\textbackslash n"} to the remote embedded device (\texttt{/dev/meascdd}), then exits the application
        \item Automatically sends STOP command first if streaming is active, ensuring clean server-side termination before shutdown
        \item Performs complete cleanup: stops network thread with \texttt{pthread\_join()}, closes socket, clears plot state, and calls \texttt{gtk\_main\_quit()}
        \item Can only be executed from CONNECTED state (not RUNNING); prevents accidental shutdown during active data acquisition
        \item Example: \texttt{SHUTDOWN} \\
        \(\rightarrow\) Dialog: ``Are you sure you want to shutdown the device \texttt{/dev/meascdd}?'' \\
        \(\rightarrow\) User clicks ``Yes'' \(\rightarrow\) Device shuts down, application exits
    \end{itemize}

\end{enumerate}

\textbf{Command Explaination (Detailed)}
\begin{enumerate}
\item The client will send ASCII text commands through the TCP socket connection, which are parsed and executed by the GUI application. Most client commands will be processed based on the command keyword, though some may require additional parameters.
\item Before processing any command, the application will first validate the command syntax and current application state. Only if the command is valid and the state permits execution will it proceed further.
\end{enumerate}

Here are the different commands a client may send:

\begin{enumerate}
    \item \texttt{START} -- Requests data streaming from server.
    \item \texttt{STOP} -- Requests to stop data streaming.
    \item \texttt{CONFIGURE} -- Requests to change a sensor's sampling frequency.
    \item \texttt{DISCONNECT} -- Requests to disconnect from server.
    \item \texttt{SHUTDOWN} -- Requests to shut down the remote device.
    \item \texttt{CONNECT} -- Requests to establish connection with server.
\end{enumerate}

\begin{itemize}
    \item Among these, three commands (\texttt{CONNECT}, \texttt{DISCONNECT},
    \texttt{SHUTDOWN}, \texttt{START}, and \texttt{STOP}) do not require any additional
    parameters. The remaining command (\texttt{CONFIGURE}) requires parameters to
    complete the client's request.
\end{itemize}

\textbf{Client Requesting Data:}
\begin{itemize}
\item In this case, the client is requesting data streaming, specifically the continuous
sensor data from the measurement device. To process this request, the client sends
the ASCII command \texttt{"START\textbackslash n"} to indicate that it wants to begin
receiving real-time sensor data.

\item Once the START command is sent, the server will respond by first transmitting a RATES
message containing the current sampling frequencies for all sensors. This ensures the
client knows the configuration before data streaming begins.

\item After sending the RATES configuration, the server will continuously transmit binary
sensor data batches. Each batch contains multiple \texttt{sensor\_data\_t} structures
(16 bytes each) with \texttt{sensor\_id}, \texttt{sensor\_value}, and \texttt{timestamp}
fields. The network receive thread processes these batches using the \texttt{recv\_all()}
function, which ensures complete message reception by looping until all expected bytes
are received.

\item The received sensor data is inserted into sensor-specific circular buffers protected
by a pthread mutex (\texttt{sample\_lock}). This prevents multiple threads from
accessing the data simultaneously, ensuring data integrity. The GUI continuously
redraws the graph at 33\,ms intervals to display the real-time sensor traces.

\item Here is the example how the client sends the command:
\end{itemize}

\begin{minted}[bgcolor = lightgray, tabsize = 4, breaklines]{c}
// Send START command via TCP socket
send(sock_fd, "START\n", 6, 0);
// Update application state
state = STATE_RUNNING;
\end{minted}

\textbf{Client Requesting to Disable a Sensor:}
\begin{itemize}
\item Similar to a configuration request, the client will disable a sensor through the GUI
by unchecking the sensor's checkbox. This is a client-side operation that does not
send any command to the server.

\item Upon unchecking a checkbox, the application will first validate the 2-sensor maximum
rule by checking whether the total selected sensors exceeds the limit. It will ensure
that users can only display up to 2 sensors simultaneously for optimal visualization.

\item If the checkbox state change is valid, the application will proceed to update the
configuration dropdown and redraw the graph area, effectively removing the sensor's
trace from the display. However, the sensor data continues to be received and stored
in the background circular buffers.

\item Here is the example of how checkbox state affects sensor display:
\end{itemize}
\begin{minted}[bgcolor = lightgray, tabsize = 4, breaklines]{c}
// When user unchecks Temperature checkbox
gtk_toggle_button_set_active(GTK_TOGGLE_BUTTON(checkboxes[0]), FALSE);

// Application updates dropdown
update_dropdown();

// Graph redraws without Temperature trace
gtk_widget_queue_draw(graph_area);
\end{minted}

\textbf{Client Requesting to Enable a Sensor:}
\begin{itemize}
\item Similar to the disable request, the client will enable a sensor by checking the sensor's checkbox. This parameter will specify which sensor the client wants to display on the graph.
\item On the client side, the process will follow the same steps as the disable request. The application will first validate the 2-sensor rule, ensuring that enabling this sensor doesn't exceed the maximum allowed sensors. If the validation passes, the application will proceed with enabling the sensor display.
\item Here is the example of checkbox enabling:
\end{itemize}
\begin{minted}[bgcolor = lightgray, tabsize = 4, breaklines]{c}
// When user checks ADC0 checkbox  
gtk_toggle_button_set_active(GTK_TOGGLE_BUTTON(checkboxes[1]), TRUE);

// Application validates 2-sensor rule and updates display
update_dropdown();
\end{minted}

\textbf{Client Requesting to Change Sample Time of a Sensor:}
\begin{itemize}
\item In this request, the client must send two parameters along with the command. Unlike previous display-only requests, this request requires communication with the server. The client must specify which sensor's sample frequency needs to be changed and the new desired sampling frequency.
\item An important condition to note is that the new sample frequency must be within the range of 10-1000 Hz, which is the valid operating range predefined by the application.
\item On the server side, the process begins with validating the sensor ID, ensuring it falls within the acceptable range (0-4 for the five sensors). In addition to this, the server will also verify the new sample frequency. If the provided sample time is less than 10 Hz or greater than 1000 Hz, the server will reject the request and send an error message back to the client. However, if both conditions are met, the server will proceed with updating the sensor's sampling timer accordingly.

\item Here is the example of client command structure:
\end{itemize}

\begin{minted}[bgcolor = lightgray, tabsize = 4, breaklines]{c}
// Send CONFIGURE command via command line
send(sock_fd, "CONFIGURE TEMP 200\n", strlen("CONFIGURE TEMP 200\n"), 0);

// Or via GUI Configure button
char net_cmd[64];
snprintf(net_cmd, sizeof(net_cmd), "CONFIGURE %s %d\n", "ADC0", 500);
send(sock_fd, net_cmd, strlen(net_cmd), 0);

// Update local frequency table
g_hash_table_replace(sensor_freq, g_strdup("TEMP"), g_strdup("200"));
\end{minted}

\textbf{Client Requesting for Disconnection:}
\begin{itemize}
\item Similar to the START command, the client will send a DISCONNECT command without any parameters. However, instead of requesting data, this command is meant to terminate the connection between the client and the server.
\item Upon receiving this request, the client application will stop the network receive thread and proceed to close the TCP socket connection, effectively disconnecting from the server. The application performs thread cleanup using \texttt{pthread\_join()} to ensure the network thread fully terminates before closing the socket. Once the disconnection is complete, the client will no longer be able to communicate with the server unless a new connection is established.
\item Here is the example of client disconnect implementation:
\end{itemize}
\begin{minted}[bgcolor = lightgray, tabsize = 4, breaklines]{c}
// Stop network thread
if (net_running)
{
      net_running = 0;                                          // Signal thread to exit
      shutdown(sock_fd, SHUT_RDWR);          // Unblock recv() call
      pthread_join(net_thread, NULL);            // Wait for thread termination
}
// Close socket
if (sock_fd >= 0)
{
      close(sock_fd);
      sock_fd = -1;
}
// Reset state
    state = STATE_DISCONNECTED;
\end{minted}

\textbf{Client Requesting for Shutdown:}
\begin{itemize}
\item Similar to other requests without parameters, this command is sent by the client without any additional data. However, unlike a simple disconnection request, this command instructs the server to shut down completely, terminating both the server application and the client GUI.
\item Once the server receives this request, it will first send STOP command if streaming is active, ensuring proper cleanup of sensor data acquisition. After stopping data collection, the server will then proceed to terminate the network thread using \texttt{pthread\_join()}, ensuring that all running threads are properly closed. This step is crucial to prevent any resources from being left open or mismanaged. Once all threads have been successfully terminated, both the client and server will complete the shutdown process and exit the application using \texttt{gtk\_main\_quit()}.
\item Here is the example of client shutdown sequence:
\end{itemize}
\begin{minted}[bgcolor = lightgray, tabsize = 4, breaklines]{c}
// If streaming, stop first
if (state == STATE_RUNNING && sock_fd >= 0)
{
    send(sock_fd, "STOP\n", 5, 0);
}

// Send shutdown command
send(sock_fd, "SHUTDOWN\n", 9, 0);

// Cleanup network resources
if (net_running)
{
    net_running = 0;
    shutdown(sock_fd, SHUT_RDWR);
    pthread_join(net_thread, NULL);
}

// Close socket and exit
close(sock_fd);
gtk_main_quit();
\end{minted}